\documentclass[11pt,a4paper]{article}

\usepackage[utf8]{inputenc}
\usepackage[czech]{babel}
\usepackage[T1]{fontenc}
\usepackage{lmodern}
\usepackage{graphicx}
\usepackage{float}
\usepackage{hyperref}
\usepackage{listings}
\usepackage{xcolor}
\usepackage[
    textwidth=165mm,
    textheight=245mm,
    centering
]{geometry}

\hypersetup{
    colorlinks=true,
    linkcolor=black,
    urlcolor=blue
}

\lstset{
    basicstyle=\ttfamily\small,
    breaklines=true,
    frame=single,
    columns=fullflexible
}

\title{IVH -- Projekt \\ Čítač na maticovém displeji}
\author{Michael Kubát (x225695)}
\date{2026}

\begin{document}

\begin{titlepage}
    \centering
    {\Huge \textsc{Vysoké učení technické v Brně}} \\
    \vspace{0.5em}
    {\huge \textsc{Fakulta informačních technologií}} \\
    \vspace{\stretch{0.35}}
    {\huge IVH -- Projekt} \\
    \vspace{0.8em}
    {\Huge Čítač na maticovém displeji} \\
    \vspace{1.5em}
    {\Large FPGA / Zynq / VHDL} \\
    \vspace{\stretch{0.65}}
    {\Large 2026 \hfill Michael Kubát (x225695)}
\end{titlepage}

\tableofcontents
\newpage

\section{Úvod}

Cílem projektu bylo navrhnout a implementovat čítač zobrazovaný na maticovém displeji 8x8. Návrh je určen pro FPGA desku se Zynq systémem. Projekt je vytvořen ve VHDL a propojen ve Vivadu pomocí blokového návrhu, který obsahuje hodinový generátor a hlavní komponentu \texttt{top}.

Výsledná aplikace postupně zobrazuje číselnou sekvenci, poté přechází do zobrazení obrázku a vlastní animace. Po dokončení celé sekvence se systém vrací zpět na začátek. Zobrazování na displeji je řešeno multiplexováním sloupců.

\section{Požadované chování systému}

Systém pracuje podle následující posloupnosti:

\begin{itemize}
    \item čítání od 00 do 10 s krokem 1 sekunda,
    \item čítání od 10 do 99 s krokem 100 ms,
    \item inverzní odpočet 90, 80, ..., 00 s krokem 500 ms,
    \item zobrazení obrázku srdce po dobu 5 sekund,
    \item dvě rotace obrázku doleva,
    \item vlastní animace tvořená zaplněním displeje, blikáním a spirálovým zhasínáním,
    \item návrat na začátek sekvence.
\end{itemize}

\section{Blokový návrh systému}

Na nejvyšší úrovni je návrh tvořen blokovým designem ve Vivadu. Vstupní hodinový signál \texttt{sysclk} je přiveden do komponenty Clocking Wizard, která generuje hodinový signál 25 MHz pro hlavní VHDL entitu \texttt{top}. Z entity \texttt{top} jsou vyvedeny výstupy \texttt{row[7:0]} a \texttt{col[7:0]} pro řízení maticového displeje.

\begin{figure}[H]
    \centering
    % Sem vlozit screenshot blokoveho designu z Vivada.
    \includegraphics[width=0.85\textwidth]{Blokovy_navrh.png}
    \caption{Blokový návrh ve Vivadu: Clocking Wizard a hlavní entita top.}
    \label{fig:block_design}
\end{figure}

Vnitřní architektura entity \texttt{top} je znázorněna na obrázku~\ref{fig:architecture}. Návrh se skládá z časovačů, FSM řadiče, dvojice BCD čítačů, dvou ROM pamětí pro převod číslic na sedmisegmentový kód, generátoru bitmapy, 64 buněk \texttt{cell} a zobrazovací části pro multiplexování displeje.

\begin{figure}[H]
    \centering
    % Sem vlozit vlastni blokove schema architektury.
    \includegraphics[width=0.95\textwidth]{architektura.png}
    \caption{Vnitřní architektura systému.}
    \label{fig:architecture}
\end{figure}

\section{Popis architektury}

\subsection{Hlavní FSM řadič}

Chování systému je řízeno konečným automatem. V návrhu jsou použity tyto stavy:

\begin{lstlisting}
S_START_RESET
S_UP_SLOW
S_UP_FAST
S_DOWN_INV
S_IMAGE
S_ROTATE
S_FILL
S_BLINK
S_SPIRAL
\end{lstlisting}

Stav \texttt{S\_START\_RESET} slouží ke krátkému počátečnímu resetu BCD čítačů. Ve stavu \texttt{S\_UP\_SLOW} probíhá pomalé čítání od 00 do 10. Ve stavu \texttt{S\_UP\_FAST} systém pokračuje rychlým čítáním do 99. Po dosažení hodnoty 99 následuje stav \texttt{S\_DOWN\_INV}, kde se zobrazuje inverzní odpočet po desítkách. Poté následuje zobrazení obrázku, jeho rotace a vlastní animace.

\subsection{BCD čítače}

Pro čítání jsou použity dvě instance komponenty \texttt{counter}. První instance reprezentuje jednotky a druhá desítky. Čítače mají vstupy pro reset, směr čítání a povolení čítání. Směr čítání je řízen signálem \texttt{dir\_cnt}. V normálních stavech systém čítá nahoru, při inverzním odpočtu čítá dolů.

Jednotky jsou v inverzním odpočtu zobrazeny jako nula, aby vznikla posloupnost:

\begin{center}
90, 80, 70, ..., 00
\end{center}

\subsection{ROM paměť pro sedmisegmentový kód}

Pro převod BCD číslice na sedmisegmentový kód je použita komponenta \texttt{seg\_rom}. Tato ROM paměť převádí hodnoty 0 až 9 na vektor segmentů \texttt{a} až \texttt{g}. V návrhu jsou použity dvě instance, jedna pro desítky a jedna pro jednotky.

Výstup ROM má tvar:

\begin{center}
\texttt{seg(6 downto 0) = a b c d e f g}
\end{center}

Hodnota 1 znamená, že daný segment má svítit.

\subsection{Generování 8x8 bitmapy}

Výstupy ROM jsou převedeny na 64bitovou bitmapu \texttt{base\_matrix}. Bitmapa reprezentuje displej 8x8 a je uložena po řádcích. Hodnota 1 znamená, že daná LED má svítit.

Číslice jsou vykresleny pomocí pomocných procedur \texttt{set\_pixel}, \texttt{hline}, \texttt{vline} a \texttt{draw\_digit}. Kvůli fyzické orientaci displeje jsou při vykreslování prohozeny horní a dolní segmenty.

\subsection{Buňky cell}

Výsledný obraz není posílán přímo do zobrazovače, ale nejprve prochází přes 64 instancí entity \texttt{cell}. Tyto instance jsou vytvořeny pomocí konstrukce \texttt{for-generate}. Každá buňka vybírá, jaký zdroj obrazu bude použit.

\newpage

\noindent Použité režimy buněk jsou:

\begin{center}
\begin{tabular}{|c|l|}
\hline
Režim & Význam \\
\hline
\texttt{000} & BCD obraz \\
\texttt{001} & BCD obraz pro inverzní odpočet \\
\texttt{010} & statický obrázek nebo rotace \\
\texttt{100} & vlastní animace \\
\hline
\end{tabular}
\end{center}

\subsection{Multiplexování displeje}

Maticový displej je řízen multiplexováním sloupců. Vždy je aktivní právě jeden sloupec. Sloupce jsou aktivní v logické 1 a řádky jsou aktivní v logické 0. Z tohoto důvodu je výstup \texttt{row} negací interní bitmapy.

Jeden sloupec je aktivní přibližně 2 ms. Po průchodu všech osmi sloupců je do zobrazovače vložen nový stabilní snímek \texttt{shown\_matrix}. Tím se omezuje problikávání při změně zobrazovaného obrazu.

\section{Použité konstanty}

V návrhu je použit generický parametr:

\begin{lstlisting}
CLK_HZ : integer := 25_000_000
\end{lstlisting}

Ten určuje frekvenci hodinového signálu. Ve skutečné implementaci je hodnota 25 MHz. V simulacích byla tato hodnota zmenšena, aby bylo možné rychleji ověřit chování systému.

\begin{center}
\begin{tabular}{|l|l|l|}
\hline
Konstanta & Hodnota & Význam \\
\hline
\texttt{TICK\_1S\_MAX} & \texttt{CLK\_HZ - 1} & časový pulz 1 s \\
\texttt{TICK\_100MS\_MAX} & \texttt{CLK\_HZ / 10 - 1} & časový pulz 100 ms \\
\texttt{TICK\_500MS\_MAX} & \texttt{CLK\_HZ / 2 - 1} & časový pulz 500 ms \\
\texttt{SCAN\_MAX} & \texttt{CLK\_HZ / 500 - 1} & 2 ms na jeden sloupec \\
\texttt{IMAGE\_5S\_MAX} & \texttt{5 * CLK\_HZ - 1} & zobrazení obrázku 5 s \\
\texttt{ROTATE\_STEP\_MAX} & \texttt{CLK\_HZ / 4 - 1} & krok rotace 250 ms \\
\texttt{FILL\_STEP\_MAX} & \texttt{CLK\_HZ / 16 - 1} & rychlost zaplnění \\
\texttt{BLINK\_STEP\_MAX} & \texttt{CLK\_HZ / 4 - 1} & rychlost blikání \\
\texttt{SPIRAL\_STEP\_MAX} & \texttt{CLK\_HZ / 32 - 1} & rychlost spirály \\
\hline
\end{tabular}
\end{center}

Použití konstant umožňuje snadno měnit časování celého systému a zároveň zrychlit simulace pomocí menší hodnoty \texttt{CLK\_HZ}.

\section{Simulace}

\subsection{Simulace zobrazování a multiplexu}

Tato simulace ověřuje správnou funkci multiplexování maticového displeje. Ve waveformu jsou sledovány zejména signály \texttt{col\_idx}, \texttt{col}, \texttt{row}, \texttt{shown\_matrix} a pomocný signál \texttt{debug\_leds}.

Signál \texttt{col\_idx} postupně prochází hodnoty 0 až 7. Podle něj je aktivní právě jeden bit výstupu \texttt{col}. Signál \texttt{debug\_leds} ukazuje neinvertovaná data aktuálního sloupce. Výstup \texttt{row} je jejich negací, protože řádky displeje jsou aktivní v logické 0.

\begin{figure}[H]
    \centering
    % Sem vlozit screenshot simulace multiplexu.
    \includegraphics[width=0.95\textwidth]{ciselny_prechod.png}
    \caption{Simulace zobrazování a multiplexování displeje.}
    \label{fig:sim_multiplex}
\end{figure}

\subsection{Simulace FSM řadiče}

Simulace FSM řadiče ověřuje průchod všemi hlavními stavy aplikace. Ve waveformu jsou sledovány zejména signály \texttt{phase}, \texttt{digit\_tens}, \texttt{digit\_ones}, \texttt{tick\_1s}, \texttt{tick\_100ms}, \texttt{tick\_500ms}, \texttt{cell\_mode} a \texttt{cell\_invert}.

\begin{figure}[H]
    \centering
    % Sem vlozit screenshot simulace FSM.
    \includegraphics[width=0.95\textwidth]{prechod_stavu.png}
    \caption{Simulace FSM řadiče.}
    \label{fig:sim_fsm}
\end{figure}

Ve stavu \texttt{S\_UP\_SLOW} jsou čítače povolovány pulzem \texttt{tick\_1s}. Po dosažení hodnoty 10 přechází systém do stavu \texttt{S\_UP\_FAST}, kde se používá rychlejší pulz \texttt{tick\_100ms}. Po dosažení hodnoty 99 následuje inverzní odpočet ve stavu \texttt{S\_DOWN\_INV}. V tomto stavu je aktivní signál \texttt{cell\_invert}. Následně systém přechází do stavů pro obrázek, rotaci a animace. FSM řadič tedy odpovídá požadovanému chování.

\subsection{Simulace ROM paměti}

Simulace ROM paměti byla provedena pomocí testbenche \texttt{seg\_rom\_tb}. Testbench postupně nastavuje vstup \texttt{digit} na hodnoty 0 až 9 a pomocí příkazů \texttt{assert} kontroluje, zda výstup \texttt{seg} odpovídá očekávanému sedmisegmentovému kódu. Dále jsou ověřeny neplatné hodnoty 10 a 15, pro které ROM vrací nulový výstup.

\begin{figure}[H]
    \centering
    % Sem vlozit screenshot simulace ROM.
    \includegraphics[width=0.95\textwidth]{ROM_pamet.png}
    \caption{Simulace ROM paměti seg\_rom.}
    \label{fig:sim_rom}
\end{figure}

Simulace proběhla bez chybových hlášení \texttt{assert}. To potvrzuje, že ROM paměť správně převádí BCD číslice na sedmisegmentový kód.

\subsection{Simulace celého systému}

Simulace celého systému byla provedena nad entitou \texttt{top}. Pro účely simulace byla snížena hodnota generického parametru \texttt{CLK\_HZ}, aby bylo možné projít větší část chování systému v rozumném čase.

Ve waveformu byly sledovány signály \texttt{phase}, \texttt{digit\_tens}, \texttt{digit\_ones}, \texttt{disp\_tens}, \texttt{disp\_ones}, \texttt{seg\_tens}, \texttt{seg\_ones}, \texttt{cell\_mode}, \texttt{cell\_invert}, \texttt{shown\_matrix}, \texttt{row} a \texttt{col}.

\begin{figure}[H]
    \centering
    % Sem vlozit prvni screenshot celeho systemu.
    \includegraphics[width=0.95\textwidth]{prechod.png}
    \caption{Simulace celého systému -- přechod mezi zobrazením srdce a vlastní animací.}
    \label{fig:sim_system_1}
\end{figure}

Na prvním průběhu je zachycen přechod mezi stavem \texttt{S\_IMAGE}, ve kterém je po dobu
5 sekund zobrazen obrázek srdce, a následujícím animačním stavem. Ve waveformu je tedy
vidět změna zobrazovaných dat z obrazového vstupu na data odpovídající animaci. Současně
lze pozorovat průběhy signálů \texttt{row} a \texttt{debug\_leds}. Signál
\texttt{debug\_leds} je vhodný zejména pro simulaci, protože přehledně ukazuje, které LED
mají v daném okamžiku logicky svítit v aktuálně zobrazovaném sloupci. Oproti tomu signál
\texttt{row} již obsahuje fyzické buzení displeje, tedy invertovanou podobu dat
(protože na maticovém displeji platí, že logická 0 znamená rozsvícení LED). Porovnání
obou průběhů proto usnadňuje kontrolu správné funkce multiplexování i ověření, že je
výstup na displej generován správně.

\begin{figure}[H]
    \centering
    % Sem vlozit druhy screenshot celeho systemu.
    \includegraphics[width=0.95\textwidth]{nova_iterace.png}
    \caption{Simulace celého systému -- konec sekvence a začátek nové iterace.}
    \label{fig:sim_system_2}
\end{figure}

Na druhém průběhu je zachycen konec celé zobrazovací sekvence a začátek nové iterace. 
Ve waveformu je vidět dokončení stavu \texttt{S\_SPIRAL}, ve kterém probíhá spirálové 
zhasínání displeje. Po dosažení poslední pozice animace se aktivuje reset čítačů a FSM 
přechází zpět do stavu \texttt{S\_UP\_SLOW}. Tím začíná nová iterace programu, ve které 
se znovu spouští čítání od hodnoty 00. Simulace potvrzuje, že systém po dokončení 
animace nezůstane stát, ale správně se vrací na začátek požadované smyčky.

\section{Implementace na kitu Xilinx Zynq }

Po syntéze, implementaci a nahrání bitstreamu na kit Xilinx Zynq byla ověřena funkce maticového displeje. Displej zobrazuje nejprve číselnou sekvenci, poté inverzní odpočet, obrázek srdce a vlastní animaci. Zobrazování je stabilní a multiplexování sloupců probíhá správně.

\section{Funkční implementace}

Video s ukázkou funkční implementace je dostupné na následujícím odkazu:

\begin{center}
    \href{https://drive.google.com/file/d/17tBV3SdgTNu4UnVacjBxkQPx3mVv7dnv/view?usp=drive_link}{\underline{Odkaz na video s funkční implementací}}
\end{center}

Před odevzdáním byl odkaz ověřen v anonymním okně prohlížeče, aby bylo zajištěno, že je video dostupné i bez přihlášení.

\section{Závěr}

V projektu byl vytvořen systém pro řízení maticového displeje 8x8. Návrh obsahuje FSM řadič, dvojici BCD čítačů, ROM paměť pro převod číslic, generování bitmapy, 64 instancí buněk \texttt{cell} a~zobrazovací multiplex. Systém splňuje požadovanou posloupnost čítání, inverzního odpočtu, zobrazení obrázku a animace.

\end{document}